% Generado por sistema/backend ReporteTex.java -- no editar a mano.
% Datos: /private/tmp/claude-501/-Users-miguel-project-taller-s/a868368b-4351-4d87-8017-fb9a5bc093aa/scratchpad/titulados.csv (titulados) / /private/tmp/claude-501/-Users-miguel-project-taller-s/a868368b-4351-4d87-8017-fb9a5bc093aa/scratchpad/empleadores.csv (empleadores)
Con 148 titulados y 53 organizaciones empleadoras, el 83.1\% de los titulados trabaja o emprende y tarda en promedio 2.63 meses en conseguir su primer empleo (IC 95\%: 1.85 a 3.42 meses); el 90.7\% considera que sus labores son pertinentes a su formación y el 74.1\% está satisfecho con ella, y lo que más se pide mejorar es ``Actualizar el plan de estudios'' (71.6\%). El 83.1\% quiere hacer un posgrado, sobre todo en ``IA'', pero solo el 41.5\% de ellos lo haría en la UMSS. Los empleadores valoran el perfil profesional (entre el 79.2\% y el 88.7\% de acuerdo), pero en la relación con la Carrera 6 de 12 aspectos tienen menos de la mitad de acuerdo. De los 14 cruces de variables, 6 muestran dependencia significativa al 5\%: ``Nivel cursado'' con ``Financiamiento'' (C7), ``Edad'' con ``Área de interés'' (C8), ``Experiencia'' con ``Remuneración'' (C9), ``Experiencia'' con ``Cargo'' (C10), ``Formación complementaria'' con ``Interés en posgrado'' (C11), ``Edad'' con ``Situación laboral'' (C12).